\section*{Chapter \ref{ch:rc:funding-margin-and-capital-adjustments} --- Funding, Margin and Capital Adjustments}

\begin{solution}{exo:rc:funding-margin-and-capital-adjustments:1}
$0.006\times5\,000\,000\times4 = \text{USD}~120\,000$.
\end{solution}

\begin{solution}{exo:rc:funding-margin-and-capital-adjustments:2}
Supervisory duration $(1-e^{-0.25})/0.05 = 4.42$; add-on $0.5\%\times50\text{ million}\times4.42 =
\text{USD}~1.106$ million; with no replacement cost, EAD $= 1.4\times1.106 = \text{USD}~1.55$ million.
\end{solution}

\begin{solution}{exo:rc:funding-margin-and-capital-adjustments:3}
The bank's \omterm{def:rc:counterparty-exposure:profiles}{expected negative exposure} exceeds its positive one: on the chapter's upward-sloping
curve, the receiver of fixed is expected to owe more than it is owed over the swap's life. The
collateral it receives on the hedge in those states is worth more (FBA) than the collateral it
posts in the others costs (FCA).
\end{solution}

\begin{solution}{exo:rc:funding-margin-and-capital-adjustments:4}
Both value the bank's own credit spread applied to what it owes: DVA through the chance of not
paying at default, FBA through funding at a rate that contains that spread. Counting both
double-counts the spread; banks either drop DVA from pricing, compute FBA on the funding spread
net of the credit spread, or use one symmetric FVA with DVA removed.
\end{solution}

\begin{solution}{exo:rc:funding-margin-and-capital-adjustments:5}
CVA charges the expected loss; KVA charges a 10\% return on capital sized for losses far out in
the tail (counterparty capital and CVA capital), held for the swap's whole life. For a BBB name,
capital is several times the expected loss.
\end{solution}

\begin{solution}{exo:rc:funding-margin-and-capital-adjustments:6}
Initial margin is required by regulation (uncleared margin rules) or by the clearing house,
must be segregated and cannot be reused; it earns less than the bank pays to fund it. It can be
reduced by netting margin across trades (portfolio models), by clearing, or by compression.
\end{solution}

\begin{solution}{exo:rc:funding-margin-and-capital-adjustments:7}
Without CSA USD~63\,700 at 40 basis points and 191\,101 at 120; with the CSA USD~1\,551 and 4\,654.
The FCA is linear in the spread for a flat spread.
\end{solution}

\begin{solution}{exo:rc:funding-margin-and-capital-adjustments:8}
The desk's hedge creates the funding need, so the cost belongs to the trade; if the desk
prices at the overnight rate, treasury subsidises uncollateralised trades and the bank
overtrades them. A funds transfer price that charges each trade its funding use, net of what
it releases, puts the cost where the decision is taken.
\end{solution}

\begin{solution}{pb:rc:funding-margin-and-capital-adjustments:1}
\textbf{1.} FCA USD~127\,400, FBA 185\,797; net FVA a benefit of 58\,397.
\textbf{2.} 1.5 basis points a year.
\textbf{3.} USD~3\,102, 98\% removed.
\textbf{4.} The client posts nothing; the hedge is collateralised, so the bank's cash needs follow
the hedge's margin, which mirrors the client trade's value.
\textbf{5.} FCA only: USD~127\,400.
\textbf{6.} CVA USD~272\,251 and DVA 210\,351 without the CSA; 33\,814 and 17\,745 with it.
\textbf{7.} USD~5.51 million without, 2.35 million with.
\textbf{8.} USD~558\,718 without, 156\,573 with.
\textbf{9.} Capital keeps a floor under margin: the add-on scaled by the margin period of risk and
the threshold-plus-MTA replacement cost, and CVA capital is charged on that EAD. The CSA cuts
KVA by 72\% against 88\% for CVA.
\textbf{10.} MVA of USD~175\,217, the funding of USD~4.22 million of initial margin falling over
the life.
\textbf{11.} USD~562\,220 without; 347\,595 with CSA and initial margin.
\textbf{12.} 6.8 and 4.2 basis points a year.
\textbf{13.} Up to the difference, 2.6 basis points a year (4.7 if no initial margin were
required, with a total of USD~172\,378, 2.1 basis points).
\textbf{14.} Daily variation margin in cash or eligible bonds, and possibly initial margin; the
cost is its own funding of that collateral and the operations to meet calls.
\textbf{15.} For this pair yes: a CSA without initial margin removes most credit, funding and
capital costs without MVA. Rules requiring initial margin apply to large financial
counterparties, generally not to corporates.
\textbf{16.} The shareholder view for the desk's quote (it pays the funding), the firm-value view
for the fair value the accounts report; hence the conventions of chapter~20.
\textbf{17.} Because capital is expensive equity that must earn its \omterm{def:rc:funding-margin-and-capital-adjustments:kva}{hurdle rate}; a trade that does
not pay for its capital destroys value for shareholders even if it is profitable before capital.
\textbf{18.} Through the funding spread (FCA, FBA, MVA), the bank's survival in first-to-default
weights, and DVA.
\textbf{19.} Named result: \emph{the funding charge}: FCA of USD~127\,400 on the uncollateralised
ten-year swap at 80 basis points, of which the CSA removes 98\%.
\textbf{20.} Lower credit, funding and capital charges, paid for with operations, liquidity and
margin costs.
\end{solution}

\begin{solution}{iq:rc:funding-margin-and-capital-adjustments:1}
The bank hedges the client's uncollateralised swap with a collateralised one. When the client
owes the bank, the bank posts margin on the hedge and borrows it at its funding rate; when the
bank owes the client, it receives margin on the hedge. FVA is the expected cost of the first
less the benefit of the second.
\iqlookfor{the hedge's margin as the source of funding.}
\end{solution}

\begin{solution}{iq:rc:funding-margin-and-capital-adjustments:2}
Firm-value view: funding costs reflect the bank's default risk, a transfer to creditors, so the
value of the firm does not change and prices should not include FVA. Shareholder view: that
transfer is a cost to shareholders, so a desk maximising shareholder value must charge it. A
defensible answer separates fair value from the desk's pricing.
\iqlookfor{both views and a clear position.}
\end{solution}

\begin{solution}{iq:rc:funding-margin-and-capital-adjustments:3}
Replacement cost $\max(V-C,0)$; add-on $0.5\%\times$ notional $\times$ supervisory duration
$\times$ maturity factor $\times$ delta, aggregated across the maturity buckets of its currency;
multiplier 1 unless over-collateralised or out of the money; EAD $= 1.4(\mathrm{RC}+\mathrm{PFE})$.
\iqlookfor{RC, add-on ingredients and alpha.}
\end{solution}

\begin{solution}{iq:rc:funding-margin-and-capital-adjustments:4}
The cost of funding segregated initial margin over the trade's life, at the funding spread over
what margin earns. Reduce it by portfolio margining, trading with counterparties whose risk
offsets, clearing, compression, and posting securities cheaper to fund than cash.
\iqlookfor{the definition and the levers.}
\end{solution}

\begin{solution}{iq:rc:funding-margin-and-capital-adjustments:5}
It depends on regulation that may change, on the \omterm{def:rc:funding-margin-and-capital-adjustments:kva}{hurdle rate} (a management choice), and on
projecting capital decades ahead; it can double-count with other charges. For a new trade:
project counterparty, CVA and market-risk capital with and without it along scenarios, take the
difference, and charge the \omterm{def:rc:funding-margin-and-capital-adjustments:kva}{hurdle rate} on its discounted path.
\iqlookfor{the doubts and the incremental method.}
\end{solution}

\begin{solution}{iq:rc:funding-margin-and-capital-adjustments:6}
Common scenarios and one exposure engine for all \omterm{def:rc:counterparty-exposure:netting}{netting sets}; profiles stored once and reused by
CVA, DVA, FCA, FBA and KVA; regulatory capital computed from the same trade representation;
algorithmic differentiation for sensitivities; parallel over paths and \omterm{def:rc:counterparty-exposure:netting}{netting sets}; incremental
runs for new trades.
\iqlookfor{shared scenarios, reuse and AAD.}
\end{solution}
